\documentclass[11pt,letterpaper]{article}
\usepackage[T1]{fontenc}
\usepackage{lmodern}
\usepackage[margin=28mm]{geometry}
\usepackage{amsmath,amssymb,amsthm}
\usepackage{microtype}
\usepackage[hidelinks]{hyperref}

\newtheorem{theorem}{Theorem}
\newtheorem{lemma}[theorem]{Lemma}

\newcommand{\R}{\mathbb R}
\newcommand{\C}{\mathbb C}
\newcommand{\Sym}{\mathcal S}
\newcommand{\Skew}{\mathcal K}
\newcommand{\ip}[2]{\langle #1,#2\rangle_F}
\newcommand{\diag}{\operatorname{diag}}
\newcommand{\rank}{\operatorname{rank}}
\newcommand{\doi}[1]{\mbox{doi \href{https://doi.org/#1}{\nolinkurl{#1}}}}

\title{Symmetric maximizers of Lyapunov operators}
\author{Shanmu Jin\thanks{Department of Neurosurgery, Peking Union Medical College
  Hospital, No.~1 Shuaifuyuan, Dongcheng District, Beijing
  100730, China.
  \href{mailto:jinshanmu1129@pku.edu.cn}{\nolinkurl{jinshanmu1129@pku.edu.cn}}.}}
\date{}
\hypersetup{pdftitle={Symmetric maximizers of Lyapunov operators},
  pdfauthor={Shanmu Jin}}

\begin{document}
\maketitle

\begin{abstract}
The Frobenius-induced norm of \(X\mapsto AX+XA^T\) is attained at a
symmetric matrix for every real \(n\times n\) matrix \(A\) exactly
when \(n\le 6\). The proof in order six uses a Chern number to choose
a three-dimensional subspace adapted to a leading skew-symmetric
singular pair. In the resulting block form, a symmetric pair gives
the required norm bound. An exact order-seven example and zero
padding yield counterexamples in every order \(n\ge7\), including
integer examples with every entry nonzero.
\end{abstract}

\section{Introduction}

For \(A\in\R^{n\times n}\), let \(L_A(X)=AX+XA^T\). Write
\(\|\cdot\|_F\) for the norm from the Frobenius inner product
\(\ip{X}{Y}=\operatorname{tr}(X^*Y)\) and \(\|L_A\|\) for the
induced operator norm. The symmetric
and skew-symmetric subspaces \(\Sym_n\) and \(\Skew_n\) are
orthogonal and invariant under \(L_A\). Set
\[
s(A)=\|L_A|_{\Sym_n}\|,\qquad k(A)=\|L_A|_{\Skew_n}\|,
\]
where \(k(A)=0\) if \(n=1\). Thus \(\|L_A\|=\max\{s(A),k(A)\}\),
and a symmetric matrix attains \(\|L_A\|\) exactly when \(k(A)\le s(A)\).

The question goes back to Byers and Nash~\cite{byers-nash}. Chen and
Tian formulated the inequality as a conjecture and proved it
through order five~\cite{chen-tian}. Kressner and Vandereycken gave
counterexamples in all orders from seven onward, leaving order six
open~\cite{kressner-vandereycken}.

\begin{theorem}\label{thm:main}
For a positive integer \(n\), the inequality \(k(A)\le s(A)\)
holds for every \(A\in\R^{n\times n}\) if and only if \(n\le6\).
For each \(n\ge7\), there is a counterexample whose entries are
all nonzero integers.
\end{theorem}

\section{Selecting a three-dimensional subspace}\label{sec:selection}

\begin{lemma}\label{lem:selection}
Let \(W\) be a Hermitian space with an algebraic direct sum
decomposition \(W=E_1\oplus E_2\oplus E_3\), where
\(\dim_{\C}E_i=2\). Let \(B:W\to W\) be complex linear and
let \(\alpha,\beta\) be alternating complex bilinear forms on \(W\).
For some choice of lines \(\ell_i\subset E_i\), their sum
\(L=\bigoplus_i\ell_i\) satisfies
\[
\rank g_L\le2,\qquad
\alpha|_{L\cap(\ker g_L)^\perp}=0,\qquad
\beta|_{\operatorname{im}g_L}=0,
\]
where \(g_L=P_{L^\perp}B|_L\) and \(P_{L^\perp}\) denotes orthogonal
projection.
\end{lemma}

\begin{proof}
On \(M=\prod_{i=1}^3\mathbb P(E_i)\), let \(\mathcal L_i\) be the
tautological line bundles, with fibers \(\ell_i\). The bundle
\(\mathcal L=\bigoplus_i\mathcal L_i\) has fiber \(L\), and
\(\mathcal Q=(M\times W)/\mathcal L\) has fiber \(W/L\), which we
identify with \(L^\perp\) by orthogonal projection.

Consider \(\pi:Z=\mathbb P(\mathcal L)\to M\), the bundle of lines
in the fibers of \(\mathcal L\). Write \(\mathcal T\) for its
tautological line bundle and \(\mathcal H=\mathcal T^\perp\) for its
orthogonal complement in \(\pi^*\mathcal L\). We use the same notation
for \(\mathcal L\) and \(\mathcal Q\) pulled back to \(Z\).
The manifold \(Z\) is compact of complex dimension five, and
\[
\mathcal F=(\mathcal T^*\otimes\mathcal Q)
    \oplus (\det\mathcal H)^* \oplus (\det\mathcal H)^*
\]
is a rank-five complex bundle. It has a smooth section whose value
at \((L,\ell)\), with \(H=L\cap\ell^\perp\), is
\[
\bigl(g_L|_\ell,\ \alpha|_H,\ (g_L|_H)^*\beta\bigr).
\]
Here the last component is the form
\((x,y)\mapsto\beta(g_Lx,g_Ly)\) on \(H\).
At a zero, \(\ell\subseteq\ker g_L\), so \(g_L\) has rank at most two,
\(L\cap(\ker g_L)^\perp\subseteq H\), and
\(\operatorname{im}g_L=g_L(H)\). The last two components give the
required restrictions of the forms.

It suffices to show \(c_5(\mathcal F)\ne0\). A nowhere vanishing
section would split off a trivial line, and the Whitney product
formula~\cite{grothendieck} would give \(c_5(\mathcal F)=0\). Put
\[
h_i=c_1(\mathcal L_i^*),\qquad
h=h_1+h_2+h_3,\qquad e_2=\sum_{i<j}h_i h_j,\qquad e_3=h_1h_2h_3.
\]
The classes satisfy \(h_i^2=0\), \(h^2=2e_2\) and \(he_2=3e_3\),
and \(\int_M e_3=1\) in the complex orientation.
The bundle \(\mathcal Q\) splits into the quotient lines
\(E_i/\ell_i\), whose first Chern classes are \(h_i\) because
the \(E_i\)-bundles are trivial. Put \(t=c_1(\mathcal T^*)\).
The splitting \(\mathcal L=\mathcal T\oplus\mathcal H\)
gives \(c_1((\det\mathcal H)^*)=h-t\) and hence
\[
c_5(\mathcal F)=(h-t)^2\prod_{i=1}^3(t+h_i).
\]

The tautological inclusion is a nowhere vanishing section of
\(\mathcal T^*\otimes\mathcal L\), so its top Chern class is zero.
Hence \(\prod_i(t-h_i)=0\), or \(t^3-ht^2+e_2t-e_3=0\).
Fiber integration \(\pi_*\) over \(\mathbb P^2\) sends
\(1,t,t^2\) to \(0,0,1\), respectively, because \(t\) restricts to
the hyperplane class. By the projection formula, the relation gives
\(\pi_*t^3=h\) and \(\pi_*t^4=h^2-e_2=e_2\). Using
\(\prod_i(t+h_i)=2(ht^2+e_3)\), we obtain
\[
\pi_*c_5(\mathcal F)
 =2h(h^2-2h^2+e_2)+2e_3
 =-4e_3.
\]
Hence \(\int_Z c_5(\mathcal F)=-4\), and the section has a zero.
\end{proof}

\section{The comparison in order six}\label{sec:order-six}

\begin{lemma}\label{lem:unitary}
Let \(U,V\in\Skew_6\) and \(B\in\R^{6\times6}\). There exists
a unitary \(E\in\C^{6\times6}\) such that
\begin{equation}\label{eq:unitary-blocks}
E^*U\bar E=\begin{pmatrix}0&X\\-X^T&U_{22}\end{pmatrix},\quad
E^*V\bar E=\begin{pmatrix}V_{11}&Y\\-Y^T&0\end{pmatrix},\quad
E^*BE=\begin{pmatrix}B_{11}&B_{12}\\G&B_{22}\end{pmatrix},
\end{equation}
with blocks of order three and \(G=\diag(g_1,g_2,0)\),
\(g_1\ge g_2\ge0\). The leading \(r\times r\) blocks of
\(U_{22}\) and \(V_{11}\) vanish, where \(r=\rank G\).
\end{lemma}

\begin{proof}
For invertible \(U\), the matrix \(C=VU\) is self-adjoint for the
nondegenerate alternating form \(\omega_U(x,y)=x^TUy\), since
\(C^TU=UC\). Its distinct generalized eigenspaces are orthogonal
for \(\omega_U\). Indeed, a power of \(C-\lambda I\) annihilates
the space for \(\lambda\), is invertible on any other such space,
and can be moved from one argument of \(\omega_U\) to the other.
The restriction of \(\omega_U\) to each generalized eigenspace
is therefore nondegenerate, so each has even dimension. Thus the
characteristic polynomial of \(C\) is the square of a monic cubic.

Comparing coefficients in descending order expresses the cubic
coefficients as polynomials in the first three coefficients below
the leading term of the square. The cubic discriminant therefore
extends to a polynomial in the entries of \(U,V\). To see that
its product with \(\det U\det V\) is nonzero, take
\[
U=\diag(J,2J,3J),\qquad V=\diag(J,J,J),\qquad
J=\begin{pmatrix}0&1\\-1&0\end{pmatrix},
\]
where the cubic has roots \(-1,-2,-3\). Consequently the pairs
with \(U,V\) invertible and three distinct cubic roots are dense.

For such a pair, the three generalized eigenspaces have dimension
two. On a two-dimensional space a self-adjoint map for a
nondegenerate alternating form is scalar, since
\(M^TJ=JM\) forces a \(2\times2\) matrix \(M\) to be scalar.
Apply Lemma~\ref{lem:selection} to the three eigenspaces of \(C\)
with \(\alpha(x,y)=x^TVy\) and \(\beta(x,y)=x^TUy\).
The resulting \(L=\bigoplus_i\ell_i\) is isotropic for \(\omega_U\)
and satisfies \(CL=L\).
Because \(U\) is real and invertible, \(U\bar L=L^\perp\).
Indeed \(x^*U\bar y=\overline{x^TUy}=0\) for \(x,y\in L\),
and both spaces have dimension three. Taking conjugates gives
\(\overline{L^\perp}=UL\), so \(V\overline{L^\perp}=CL=L\).
These two identities give the zero diagonal blocks in
\eqref{eq:unitary-blocks} in any unitary basis adapted to \(L\).

Choose orthonormal bases of \(L\) and \(L^\perp\) using a singular
value decomposition of \(g_L\). Then \(G\) has the required diagonal
form. The first \(r\) vectors in these bases span
\(L\cap(\ker g_L)^\perp\) and \(\operatorname{im}g_L\), respectively.
The forms vanish on these spaces, and reality of \(U,V\) gives the
same zeros in the congruence blocks \(V_{11}\) and \(U_{22}\).

For arbitrary \(U,V\), approximate them by pairs from the dense set
with \(B\) fixed. Compactness of the unitary group gives a convergent
subsequence of the bases. The displayed block form passes to the
limit. Every positive singular value of the limiting \(G\) is
positive along that subsequence eventually, so the leading
\(r\times r\) zeros also pass to the limit.
\end{proof}

\begin{proof}[Proof of Theorem~\ref{thm:main} in order six]
The case \(k(A)=0\) is immediate. For \(k(A)>0\), set \(B=A/k(A)\)
and choose unit skew-symmetric matrices \(U,V\) with
\[
L_B(V)=U,\qquad L_{B^T}(U)=V.
\]
Such a pair comes from a largest singular value of \(L_B|_{\Skew_6}\),
whose adjoint is \(L_{B^T}|_{\Skew_6}\).

Complexification preserves the restricted norms, since the squared
Frobenius norms of the real and imaginary parts add.
Unitary congruence \(Z\mapsto E^*Z\bar E\) preserves the norm
and both transpose symmetry classes. It sends \(L_B\) and its
adjoint to \(L_{\widehat B}\) and \(L_{\widehat B^*}\), with
\(\widehat B=E^*BE\). Thus the skew pair \(\widehat U,\widehat V\)
in \eqref{eq:unitary-blocks} satisfies
\(L_{\widehat B}(\widehat V)=\widehat U\) and
\(L_{\widehat B^*}(\widehat U)=\widehat V\).

For \(\varepsilon\in\{1,-1\}\), define
\[
\begin{aligned}
S_\varepsilon&=XG+\varepsilon(XG)^T,&
T_\varepsilon&=GY+\varepsilon(GY)^T,\\
U_\varepsilon&=\begin{pmatrix}0&X\\\varepsilon X^T&T_\varepsilon\end{pmatrix},&
V_\varepsilon&=\begin{pmatrix}S_\varepsilon&Y\\\varepsilon Y^T&0\end{pmatrix}.
\end{aligned}
\]
The bottom-right block of \(L_{\widehat B}(\widehat V)=\widehat U\)
and the top-left block of \(L_{\widehat B^*}(\widehat U)=\widehat V\)
give \(U_{22}=T_{-1}\) and \(V_{11}=S_{-1}\), since \(G\) is real
diagonal. Thus \(U_{-1}=\widehat U\) and \(V_{-1}=\widehat V\),
with at least one of \(X,Y\) nonzero. Write
\[
\Delta_\varepsilon
 =\|U_\varepsilon\|_F^2+\|V_\varepsilon\|_F^2
  -2\operatorname{Re}\ip{U_\varepsilon}{L_{\widehat B}(V_\varepsilon)}.
\]
The singular equations give \(\Delta_{-1}=0\).
The upper-right block of \(L_{\widehat B}(V_\varepsilon)\) is
\(B_{11}Y+YB_{22}^T+S_\varepsilon G^T\), and its bottom-right
block is \(T_\varepsilon\). Transpose preserves the Frobenius
inner product, so the two off-diagonal blocks contribute equally.
Since \(Z\mapsto(Z+\varepsilon Z^T)/2\) is the orthogonal
projection onto \(Z^T=\varepsilon Z\), we have
\(\ip{X}{S_\varepsilon G^T}=\ip{XG}{S_\varepsilon}
=\|S_\varepsilon\|_F^2/2\). It follows that
\[
\Delta_\varepsilon
 =2\bigl(\|X\|_F^2+\|Y\|_F^2\bigr)
  -4\operatorname{Re}\ip{X}{B_{11}Y+YB_{22}^T}
  -\|S_\varepsilon\|_F^2-\|T_\varepsilon\|_F^2 .
\]
Subtracting the two cases gives
\[
\Delta_1-\Delta_{-1}
  =-4\bigl(\ip{XG}{(XG)^T}+\ip{GY}{(GY)^T}\bigr).
\]
If \(r=\rank G\), the entries of \(XG\) beyond the first \(r\)
columns and those of \(GY\) beyond the first \(r\) rows are zero.
In both matrices the leading \(r\times r\) block is symmetric,
because the corresponding blocks of \(S_{-1}\) and \(T_{-1}\)
vanish by Lemma~\ref{lem:unitary}. Each inner product is therefore
the squared Frobenius norm of that leading block, and
\(\Delta_1\le0\).

The matrices \(U_1,V_1\) are symmetric. By Cauchy--Schwarz and the
arithmetic--geometric mean inequality,
\[
\Delta_1\ \ge\
(1-s(B))\bigl(\|U_1\|_F^2+\|V_1\|_F^2\bigr).
\]
The sum of squared norms is positive, so \(s(B)\ge1\).
Scaling back gives \(s(A)\ge k(A)\).
\end{proof}

\section{Zero padding and sharpness}\label{sec:sharpness}

For any real \(A\) and \(m\ge1\), let \(0_m\) denote the square
zero matrix of order \(m\). In either transpose symmetry class, the
two diagonal blocks and the off-diagonal pair form orthogonal
subspaces for the splitting of \(A\oplus0_m\)~\cite{kressner-vandereycken}.
The Lyapunov
operator acts on these as \(L_A\), zero and \(Y\mapsto AY\),
respectively. The two off-diagonal blocks have the same Frobenius
norm, so this last map contributes \(\|A\|_2\), the spectral norm.
For any real unit vector \(x\),
\[
\|L_A(xx^T)\|_F^2=2\|Ax\|_2^2+2(x^TAx)^2 .
\]
Taking a leading right singular vector of \(A\) gives
\(s(A)\ge\sqrt{2}\|A\|_2\). Hence
\begin{equation}\label{eq:padding}
s(A\oplus0_m)=s(A),\qquad
k(A\oplus0_m)=\max\{k(A),\|A\|_2\}.
\end{equation}
For \(n<6\), applying the order-six inequality to \(A\oplus0_{6-n}\)
gives \(k(A)\le k(A\oplus0_{6-n})\le s(A)\).
If \(k(A)>s(A)\), both norms are preserved
under zero padding.

An order-seven matrix with a strict gap is
\begin{equation}\label{eq:example}
A_7=\begin{pmatrix}
   0 &  80 &   0 &  -1 &   0 &   0 &   7\\
 -80 & -20 & -60 &   5 &   1 &   0 &   8\\
   0 & -60 &  20 &   4 &   0 &   1 &   1\\
   0 &   7 &   4 & -80 &  20 &   0 &  10\\
   0 &  -1 &   0 & -20 & -40 &  80 & -10\\
   0 &   0 &   1 &   0 & -80 &  40 &   0\\
   7 &  -6 &   0 &  -5 &  -5 &   0 & -80
\end{pmatrix}.
\end{equation}
We verify the bounds
\begin{equation}\label{eq:bounds}
s(A_7)<\frac{1499}{8}<\frac{375}{2}<k(A_7).
\end{equation}
Let \(E_{ij}\) be the matrix units. Use \(E_{11},\dots,E_{77}\)
followed by \(E_{ij}+E_{ji}\), \(i<j\), as the symmetric basis.
Use \(E_{ij}-E_{ji}\), \(i<j\), as the skew-symmetric basis.
In both lists the pairs \((i,j)\) are ordered lexicographically.
For the relevant basis \(F_1,\dots,F_N\) and positive integers
\(p,q\), form the integer symmetric matrix
\[
M^{(1)}_{ij}=p^2\ip{F_i}{F_j}
       -q^2\ip{L_{A_7}(F_i)}{L_{A_7}(F_j)} .
\]
Its quadratic form at the coordinates of \(X\) is
\(p^2\|X\|_F^2-q^2\|L_{A_7}(X)\|_F^2\).
For a positive definite form the restricted norm is below \(p/q\).
A negative value places the norm above \(p/q\).

Use rational symmetric elimination with the recurrence
\[
d_j=M^{(j)}_{jj},\qquad
M^{(j+1)}_{i\ell}=M^{(j)}_{i\ell}
    -M^{(j)}_{ij}M^{(j)}_{j\ell}/d_j\qquad(i,\ell>j).
\]
The pivot floors obtained from \eqref{eq:example} are summarized by
\[
\begin{array}{lrrr|rr}
\text{Subspace} & N & p & q & \min_{1\le j<N}\lfloor d_j\rfloor
                            & \lfloor d_N\rfloor\\ \hline
\text{Symmetric}      &28&1499&8&130009&1024128\\
\text{Skew-symmetric} &21&375 &2& 17610&  -9319
\end{array}
\]
Each elimination is a congruence to the direct sum of the pivot
and the trailing matrix by completing the square.
All pivots in the symmetric case
are positive, and the last pivot in the skew-symmetric case is
negative. This proves \eqref{eq:bounds}.

For \(n\ge7\), set \(A_n=A_7\oplus0_{n-7}\), with the zero block
omitted at \(n=7\). Equation~\eqref{eq:padding} preserves
\eqref{eq:bounds}. Let \(\mathbf 1_n\) be the vector of ones and
set \(P_n=\mathbf 1_n\mathbf 1_n^T/(40n)\). Adding \(P_n\) to
\(A_n\) changes each restricted norm by at most \(1/20\), since
\(\|L_{P_n}\|\le2\|P_n\|_2=1/20\).
The two rational bounds in \eqref{eq:bounds} lie \(1/16\) on
either side of \(2999/16\). Since \(1/20<1/16\), the integer matrix
\[
D_n=40nA_n+\mathbf 1_n\mathbf 1_n^T
\qquad\text{satisfies}\qquad
s(D_n)<40n\cdot\frac{2999}{16}<k(D_n).
\]
Every entry of \(D_n\) is congruent to \(1\) modulo \(40n\).
This completes the proof of Theorem~\ref{thm:main}.

\section*{Acknowledgment}
This work was carried out through automated proof discovery using
GPT-6.2 Sol at the \texttt{max} reasoning setting in Codex. The human author
sorted out and reorganized the work and takes full responsibility for it.
The \LaTeX{} manuscript and Lean formalization code are available at
\url{https://github.com/jinshanmu/SymmetricMaximizer}.

\begin{thebibliography}{9}

\bibitem{byers-nash}
R.~Byers and S.~Nash.
On the singular \textquotedblleft vectors\textquotedblright\ of the
Lyapunov operator.
\emph{SIAM J. Algebraic Discrete Methods} \textbf{8}(1), 59--66 (1987).
\doi{10.1137/0608003}.

\bibitem{chen-tian}
S.~Chen and Y.~Tian.
Note on \textquotedblleft On the singular \textquoteleft vectors\textquoteright\
of the Lyapunov operator\textquotedblright\ by R.~Byers and S.~Nash.
\emph{SIAM J. Matrix Anal. Appl.} \textbf{36}(3), 1069--1072 (2015).
\doi{10.1137/140974031}.

\bibitem{kressner-vandereycken}
D.~Kressner and B.~Vandereycken.
A counterexample to the symmetric-maximizer conjecture for Lyapunov operators.
arXiv:2608.20875 (2026).
\doi{10.48550/arXiv.2608.20875}.

\bibitem{grothendieck}
A.~Grothendieck.
La th\'eorie des classes de Chern.
\emph{Bull. Soc. Math. France} \textbf{86}, 137--154 (1958).
\doi{10.24033/bsmf.1501}.

\end{thebibliography}
\end{document}
